\section*{Chapter \ref{ch:iv2:rates-and-fx} --- Rates and FX}
\markboth{\protect\omTocMark{}{Answers: Rates and FX}}{\protect\omTocMark{}{Answers: Rates and FX}}

\begin{solution}{iq:iv2:rates-and-fx:1}
It is about 82.19 with annual compounding, and about 81.87 with continuous compounding.

$100/1.04^5 = 82.19$, and $100e^{-0.2} = 81.87$. A quoted rate means nothing without its compounding basis.

\iqwrong{Using $100(1 - 5 \times 0.04) = 80$. That is a simple-discount rule, and it is 2 points off here.}
\end{solution}

\begin{solution}{iq:iv2:rates-and-fx:2}
It is about 100.96. At a yield of 5\%, it is exactly 100, as a bond with a coupon equal to its yield trades at par.

$105/1.04 = 100.96$. At 5\%, $105/1.05 = 100$.

\iqwrong{Thinking a premium means the bond is a bad deal. The coupon pays for it, and the price falls to par by maturity.}
\end{solution}

\begin{solution}{iq:iv2:rates-and-fx:3}
The Macaulay duration is 5, and the modified duration is $5/1.04 \approx 4.81$. So 10 basis points takes off about 0.481\%, which is about 40 cents on a price of 82.19.

A zero has only one payment, so its average time is its maturity. The percentage change is modified duration times the yield change.

\iqwrong{Applying the 4.81\% to \$100 of face value. It applies to the price.}
\end{solution}

\begin{solution}{iq:iv2:rates-and-fx:4}
The DV01 is \$8\,000, and you make about \$200\,000.

$10\,000\,000 \times 8 \times 0.0001 = 8\,000$, and $25 \times 8\,000 = 200\,000$. Quote DV01 in dollars, as modified duration alone does not tell you how much you own.

\iqwrong{Forgetting that a fall in yields is a gain for the owner of the bond.}
\end{solution}

\begin{solution}{iq:iv2:rates-and-fx:5}
It is about 5.01\%. To lock it in, the borrower buys a one-year zero and sells a two-year zero, which leaves the borrower receiving cash in a year and paying in two.

$1.04^2/1.03 = 1.0501$. Buying the one-year with the proceeds of selling a two-year costs nothing today. The cash flow at year 1 is an inflow, and at year 2 it is that amount grown by 5.01\%.

\iqwrong{Averaging to get 5\%. The geometric form gives 5.01\%, and over longer gaps the difference grows.}
\end{solution}

\begin{solution}{iq:iv2:rates-and-fx:6}
Short \$5.625 million of the ten-year. You are still exposed to the shape of the curve: the five-year and the ten-year may not move by the same amount.

The DV01 of your long is $10\,000\,000 \times 4.5 \times 0.0001 = \$4\,500$. The short needs that DV01: $4\,500/(8 \times 0.0001) = 5\,625\,000$.

\iqwrong{Hedging the same notional. That leaves you short \$3\,500 of DV01, and a large bet on rates.}
\end{solution}

\begin{solution}{iq:iv2:rates-and-fx:7}
You gain if rates rise. It is like being short a five-year 4\% bond, funded by a floating rate note. If the swap rate rises 10 basis points, you make about \$45\,000.

The floating leg is worth about par, and the fixed leg is a bond. Paying fixed is short that bond. So $10 \times 4\,500$.

\iqwrong{Thinking the floating leg has a lot of DV01. It resets, so it only has the duration to the next fixing.}
\end{solution}

\begin{solution}{iq:iv2:rates-and-fx:8}
The price is \$99. The yield is about 4.04\%. It is above 4\% because the discount is on the face value, while you earn your return on the lower price.

$100(1 - 0.04 \times 90/360) = 99$. The yield is $(1/99)(360/90) = 4.04\%$.

\iqwrong{Treating a discount rate as a yield. On a bill they are close, and on a long bill the gap grows.}
\end{solution}

\begin{solution}{iq:iv2:rates-and-fx:9}
The forward is about 1.1214, or about 214 points above spot. No, it does not predict that. It is set by the interest gap so that no arbitrage is possible.

$1.10 \times 1.05/1.03 = 1.12136$. The lower-rate currency trades at a forward premium. Otherwise, borrowing one and lending the other, hedged, would make money for nothing.

\iqwrong{Reading a forward above spot as a forecast of a rise. A long history of tests shows forwards to be poor forecasts.}
\end{solution}

\begin{solution}{iq:iv2:rates-and-fx:10}
EURJPY should be 165. At 166, sell EURJPY there, buy EURUSD and buy USDJPY, for 1 yen per euro. First check that all three quotes are live, what bid-ask spreads cost and when each one settles.

$1.10 \times 150 = 165$. Selling a euro for 166 yen, buying it back with 1.10 dollars, and buying those dollars for 165 yen locks in 1 yen.

\iqwrong{Using a mid-price on all three legs. Crossing three spreads usually takes away what looks like a gap this size on a screen.}
\end{solution}

\begin{solution}{iq:iv2:rates-and-fx:11}
It is a dollar put on \$1.1 million, at a strike of about 0.909 euros per dollar.

A call on 1 million euros at 1.10 gives you the right to pay \$1.1 million and to receive 1 million euros. That is the right to sell \$1.1 million at $1/1.10$ euros per dollar.

\iqwrong{Quoting a million dollars of notional. The other side's notional is the first times the strike.}
\end{solution}

\begin{solution}{iq:iv2:rates-and-fx:12}
It must be at about 1.1214 for you to break even, the same as the forward. A carry trade's risk is that the currency moves. In a year, the euro needs to rise by about 1.9\% to wipe out your 2\% gain.

Borrow 1 euro, sold for \$1.10. Lend at 5\%, to get 1.155. You owe 1.03 euros. So $1.155/1.03 = 1.1214$.

\iqwrong{Calling this free money, since the forward does not on average predict the move. The risk has a skew: large losses come in the unwinds.}
\end{solution}

\begin{solution}{iq:iv2:rates-and-fx:13}
A rise costs about 7.6\%, and a fall gains about 8.4\%. The gain is larger. People like that because you gain more than you lose for the same move.

$-8 \times 1\% + \tfrac12 \times 80 \times (1\%)^2 = -7.6\%$. The convexity term has the same sign in both directions.

\iqwrong{Thinking this is free. A curve prices convexity through lower yields on long bonds. Convexity is a position with a cost.}
\end{solution}

\begin{solution}{iq:iv2:rates-and-fx:14}
The discount factors are $D_1 = 0.97087$ and $D_2 \approx 0.92420$. The two-year zero rate is about 4.02\%, and the forward is about 5.05\%.

Par pricing at two years: $0.04D_1 + 1.04D_2 = 1$. So $D_2 = (1 - 0.04 \times 0.97087)/1.04$. The zero rate is $D_2^{-1/2} - 1$, and the forward is $D_1/D_2 - 1$.

So the zero rate is above the par yield in a rising curve, and the forward is above both.

\iqwrong{Taking the two-year par yield as the two-year zero rate. The gap, which is 2 basis points here, grows with maturity and with slope.}
\end{solution}

\begin{solution}{iq:iv2:rates-and-fx:15}
The three-year par rate is about 4.93\%. The two-year swap starting in a year is about 5.99\%. The floating leg needs no model, since it is worth $D_{\text{start}} - D_{\text{end}}$, by replication.

The discount factors are 0.97087, 0.92456 and 0.86384. The par rate is $(1 - D_3)/(D_1 + D_2 + D_3)$. The forward swap rate is $(D_1 - D_3)/(D_2 + D_3)$.

Borrow at the end and lend to the start of each period, and you get each floating payment with no risk. That holds in any rates model, in the simple case with one curve.

\iqwrong{Forecasting the floating payments and discounting. If you do it with the forwards, it gives the same answer, but it hides why no model is used.}
\end{solution}

\begin{solution}{iq:iv2:rates-and-fx:16}
The barbell holds about 41\% in the ten-year and 59\% in the two-year. Its convexity is 36.5, against 25.0 for the bullet. Rates down 100: the barbell gains 4.64 and the bullet 4.58. Rates up 100: the barbell loses 4.28 and the bullet 4.33. On a real curve, the bullet often has more yield, which pays for its lack of convexity.

The weights come from duration: 1.886, 4.452 and 8.111. So $w_{10} = (4.452 - 1.886)/(8.111 - 1.886) = 0.412$.

On a flat curve, the barbell wins whatever happens, which would be an arbitrage. So a flat curve cannot stay flat in a world with volatility, and a real one charges for convexity by giving the five-year a higher yield. That is why the trade has a carry cost.

\iqwrong{Calling the barbell free convexity. Compare carry as well, and also its exposure to a change of slope.}
\end{solution}

\begin{solution}{iq:iv2:rates-and-fx:17}
You buy about \$41.2 of the ten-year and \$58.8 of the two-year. You make about 44.5 cents per \$100 if the five-year yield is up 10 basis points. You are betting that the five-year cheapens against the wings: this is a butterfly\index{butterfly (rates)}.

With $a + b = 100$ and $1.886a + 8.111b = 445.2$, we get $b = 41.2$. The gain is $4.452$ cents per basis point on the short, for 10 basis points.

\iqwrong{Calling the trade free of risk, because of no DV01. It still has exposure to slope if the wings do not move as expected, and it has a convexity position.}
\end{solution}

\begin{solution}{iq:iv2:rates-and-fx:18}
Short \$10 million of the ten-year, whose DV01 is \$8\,000, and buy about \$42.1 million of the two-year. If the two-year falls 5 and the ten-year rises 5, you make \$80\,000. If both rise 10, you make nothing.

$8\,000/190 = 42.1$. Each leg makes $8\,000 \times 5$.

A steepener gains when the slope rises, and it has no duration. It still has carry: being long the two-year and short the ten-year usually costs carry if the curve slopes up.

\iqwrong{Using equal notional. That is a huge short duration position, and it has little to do with the slope.}
\end{solution}

\begin{solution}{iq:iv2:rates-and-fx:19}
The regression says about \$5\,250 of two-year DV01, not \$8\,000. The residual volatility is about 53.6 basis points a year, against about 60.2 with the DV01-neutral hedge.

Regressing the ten-year change on the two-year gives $\beta = 0.7 \times 75/80 = 0.656$. So $8\,000 \times 0.656 = 5\,250$. The residual is $75\sqrt{1 - 0.49} = 53.6$. The DV01-neutral spread has
\[
\sqrt{80^2 + 75^2 - 2 \times 0.7 \times 80 \times 75} = 60.2.
\]

That is about \$27.6 million of two-year, at \$190 per million, as against \$42.1 million. But the beta shifts with the regime, and regressions on levels of vol can fool you. So watch it and refresh it.

\iqwrong{Thinking DV01 is the only proper hedge. It removes a parallel shift only, and real shifts are not parallel.}
\end{solution}

\begin{solution}{iq:iv2:rates-and-fx:20}
It is worth about \$1.91 per \$100, both ways.

As a bond, $4/1.03 + 104/1.03^2 = 101.91$, less 100 for the floating leg. As a gap: $(4 - 3) \times (1/1.03 + 1/1.03^2) = 1.913$.

\iqwrong{Using the 1\% gap times two years, or \$2. The two payments are discounted.}
\end{solution}

\begin{solution}{iq:iv2:rates-and-fx:21}
It is about 13 basis points, and the future's rate is above the forward's.

The rule is $\tfrac12\sigma^2T_1T_2 = 0.5 \times 0.0001 \times 5 \times 5.25 = 13.1$ basis points. A future is marked to market. Its holder gains when rates fall, and reinvests at a low rate; the holder loses when rates rise, and funds at a high rate. Since it is a disadvantage to the long, the futures price is lower, and that means its rate is higher.

\iqwrong{Ignoring it at the far end. It grows with the square of maturity, and in five years it is larger than the bid-ask spread.}
\end{solution}

\begin{solution}{iq:iv2:rates-and-fx:22}
The third bond is cheapest. Its gross basis is the lowest, at 0.6, against 1.0 and 0.8.

The cost of delivering is the bond's price less the invoice price: $93 - 92 = 1.0$, $110.8 - 110 = 0.8$ and $125.6 - 125 = 0.6$.

A factor is the price at a set yield. Away from that yield, the bonds' prices do not move in proportion. When yields are high, the long-duration bonds fall more, and become cheapest. When low, the short-duration bonds do.

\iqwrong{Choosing the lowest-priced bond. A bond's price has to be compared with the invoice amount.}
\end{solution}

\begin{solution}{iq:iv2:rates-and-fx:23}
Carry is about 1\%, and roll-down is about 0.72\%, for a total of about 1.72\%. The four-year yield can rise about 48 basis points from 3.8\%, to about 4.28\%, before you lose money.

Carry is the yield less the financing: $4 - 3$. Over the year, the bond becomes a four-year priced at 3.8\%, rather than 4\%, so its price gains $0.2 \times 3.6 = 0.72$. The break-even is $1.72/3.6 = 0.48$.

\iqwrong{Forgetting the roll. On a steep curve, it can be larger than the carry. And it only appears if the curve does not move with time.}
\end{solution}

\begin{solution}{iq:iv2:rates-and-fx:24}
The mean is about 4.26\%, and the standard deviation about 0.93\%. The chance of a negative rate is about 0.0002\%. No, it does not make it useless. It is often only a small part of the model's prices.

The mean is $5 + (3 - 5)e^{-1}$. The variance is $\sigma^2(1 - e^{-2\kappa T})/(2\kappa)$. The chance is $\Phi(-4.26/0.93)$.

A normal model has its merits: closed forms, and a volatility that matches how rates tend to move in basis points. Negative rates have occurred in some markets, so that part is not always a flaw.

\iqwrong{Rejecting a model because of a feature with no effect on the price in question. Test the effect first.}
\end{solution}

\begin{solution}{iq:iv2:rates-and-fx:25}
It is about 16.7 basis points. Long yields fall since a random rate makes the discount factor convex in it.

The integrated rate has a variance of $\sigma^2T^3/3$. So $E[e^{-\int r}] = e^{-f T + \sigma^2T^3/6}$. The yield is $f - \sigma^2T^2/6$. For $T = 10$, $\sigma^2 = 0.0001$: $0.0001 \times 100/6 = 0.00167$. This model's forward curve is $f - \sigma^2T^2/2$; a truly flat forward curve would give a zero yield of exactly its level.

That is the source of the convexity value that long bonds carry. So yields tend to fall below expected rates at long maturity.

\iqwrong{Thinking that the volatility affects only options. The bond is an average of convex functions of the rate, so it does too.}
\end{solution}

\begin{solution}{iq:iv2:rates-and-fx:26}
The caplet is worth about 7.57 basis points of notional. The same volatility is about 80 basis points, and Bachelier gives about 7.58. At a rate of 2\%, it would be about 40 basis points.

Black's formula on the forward rate, with the accrual of 0.25 and the discount factor 0.95. At the money, it is $0.95 \times 0.25 \times 0.04 \times (2\Phi(0.1) - 1) = 0.000757$. The normal form is $0.95 \times 0.25 \times 0.008 \times 0.3989$. To convert, use $\sigma_N \approx F\sigma_L$.

\iqwrong{Comparing lognormal volatility across rate levels. A move from 20\% to 40\% as the rate halves can be the same normal volatility. Traders quote normal.}
\end{solution}

\begin{solution}{iq:iv2:rates-and-fx:27}
It is about 1.44\% of notional. The annuity multiplies it since each basis point of swap rate is paid across all the future dates. One basis point of volatility is worth about 1.8 basis points of notional.

The rule is $A\sigma_N\sqrt{T/(2\pi)} = 4.5 \times 0.008 \times 0.3989$. The payoff is $A(S - K)^+$ at expiry. Using the annuity as numeraire, the swap rate is a martingale. For the sensitivity, $4.5 \times 0.0001 \times 0.3989 = 0.018\%$.

\iqwrong{Using a discount factor in place of the annuity. It treats the swap as having a single payment.}
\end{solution}

\begin{solution}{iq:iv2:rates-and-fx:28}
It is $-4.5\%$ of notional. It is a forward-starting swap at 5\%, and no model is needed.

At expiry, the two options pay $A(S - K)^+ - A(K - S)^+ = A(S - K)$. So today, in annuity units, it is $A(F - K) = 4.5 \times (4 - 5)$.

\iqwrong{Quoting a volatility for this. Volatility does not affect the difference.}
\end{solution}

\begin{solution}{iq:iv2:rates-and-fx:29}
It implies about 5.29\%, so about 29 basis points above the cash rate. So lending dollars through the swap market pays more than through cash. That is a basis, and it does not close because arbitrage needs balance sheet and credit, and the cost of those is real.

The implied rate is $1.1245/1.10 \times 1.03 - 1 = 5.294\%$.

So the trade is to borrow cash dollars at 5\%, sell them for euros, invest at 3\% and buy dollars forward. That locks in about 29 basis points. Covered interest parity was once seen as close to a law, and since the 2008 crisis it has been systematically broken. For a bank, the constraint is its capital rules and the cost of its own balance sheet.

\iqwrong{Treating this as free money. Who is able to do it, at what funding cost, and with what balance sheet, are the real questions.}
\end{solution}

\begin{solution}{iq:iv2:rates-and-fx:30}
The forward is about 1.1222. The call is worth about \$0.0536, and the put about \$0.0324. The spot delta is 0.581, and the forward delta 0.599.

The Garman-Kohlhagen\index{Garman-Kohlhagen} form treats the foreign rate as a dividend yield. $F = Se^{(r_d - r_f)T}$. So $C = e^{-r_d}[F\Phi(d_1) - K\Phi(d_2)]$, with $d_1 = (\ln(F/K) + 0.005)/0.1$.

The spot delta is $e^{-r_f}\Phi(d_1)$. It is the amount of the euro you hold today. The forward delta, $\Phi(d_1)$, is what you hold in forwards.

Parity: $C - P = e^{-r_d}(F - K) = 0.0211$.

\iqwrong{Using the stock formula with a single rate. A foreign currency pays interest, so its forward is not at the spot grown by the dollar rate alone.}
\end{solution}

\begin{solution}{iq:iv2:rates-and-fx:31}
It is about 4.87 euro cents per euro. That is 4.87\% of the euro notional and 4.87\% of the dollar notional too, since the strike and spot are the same here. Priced from the dollar side, it must use the same 10\% volatility, since the log of a rate and of its inverse have the same variance.

$0.0536/1.10 = 0.0487$. As a percentage of the dollar notional, it is $0.0536/K$. With $K = S$, the two are equal. The dollar put on $X = 1/S$ at $1/K$, per dollar, is a European put at $F' = 1/F$. The rates swap roles. It comes to about 0.0443 euros per dollar, on \$1.10: 0.0487 euros, or 0.0536 dollars in value.

\iqwrong{Thinking a USD put on EURUSD has a different smile. It's the same option. The convention, and not the risk, is what differs.}
\end{solution}

\begin{solution}{iq:iv2:rates-and-fx:32}
The call is 9.5\%, and the put is 11.5\%. The strikes are about 1.199 and 1.048.

The call volatility is at-the-money + fly + half the risk reversal, and the put's is at-the-money + fly $-$ half the risk reversal. To find the strike, solve $e^{-r_f}\Phi(d_1) = 0.25$ for the call. So $\Phi(d_1) = 0.2576$, $d_1 = -0.651$, and $K = Fe^{-d_1\sigma + \sigma^2/2} = 1.1222e^{0.0618 + 0.0045}$.

\iqwrong{Reading the strikes as fixed. In FX they are not quoted directly, and a delta is a moving strike that depends on the volatility itself.}
\end{solution}

\begin{solution}{iq:iv2:rates-and-fx:33}
It is about 16.5\% with a correlation of 0.5, and about 9.5\% with a correlation of $-0.5$.

EURJPY is the product of the other two, so $\sigma^2 = \sigma_1^2 + \sigma_2^2 + 2\rho\sigma_1\sigma_2$. So $0.01 + 0.0081 + 0.009 = 0.0271$, whose root is 16.46\%.

In reverse, three option-implied volatilities give you an implied correlation. Traders compare that with what is realised. A pair with no liquid options can be priced from the other two, and the mispricing is a correlation trade.

\iqwrong{Using a minus sign. It is a product, not a difference. For a ratio such as EURGBP, from EURUSD and GBPUSD, the sign reverses.}
\end{solution}

\begin{solution}{iq:iv2:rates-and-fx:34}
The compounded rate is about 4.144\%, and the simple average is about 4.143\%. It is known only at the end, and that means a caplet on it keeps moving during the accrual period, so its volatility does not stop at the start.

Each day multiplies by $1 + r_id_i/360$. The product is 1.0008058, and the rate is that less 1, times 360/7.

\begin{dated}{2026-11}{Overnight rates}\label{dat:iv2:rates-and-fx:ois}
In US dollars, overnight swaps are based on SOFR, with annual payments, actual/360 on both legs, and floating compounded in arrears. Conventions like these change, so check them before you trade. See the chapter's ledger.
\end{dated}

So an in-arrears caplet has a total variance that includes the accrual period, with a variance that decays through it. A rule of thumb gives one-third of the accrual period as extra variance time.

\iqwrong{Thinking a compounded rate is close enough to an average for all uses. Over longer periods and with higher rates, it is not.}
\end{solution}

\begin{solution}{iq:iv2:rates-and-fx:35}
The five-year zero is about 0.8083, a yield of about 4.26\% continuously compounded. A simulation gives 0.8082. At a very long maturity, the yield approaches $\theta - \sigma^2/(2\kappa^2)$, a level of 4.98\%.

Write $P = A(T)e^{-B(T)r}$. Substituting gives $B' + \kappa B = 1$, so $B = (1 - e^{-\kappa T})/\kappa$. And $\ln A = (\theta - \sigma^2/(2\kappa^2))(B - T) - \sigma^2B^2/(4\kappa)$.

For the check, simulate the rate, sum it by the trapezoid rule, and average the discount.

So the long yield is below $\theta$, by the same convexity value as in Ho-Lee, but limited by the mean reversion. The speed is what keeps it finite: with no reversion, the gap grows without limit.

\iqwrong{Assuming the yield approaches the mean of the rate. The convexity term lowers it, a little here.}
\end{solution}

\begin{solution}{iq:iv2:rates-and-fx:36}
The drift is $\theta(t) = \partial_tf(0, t) + \kappa f(0, t) + \sigma^2(1 - e^{-2\kappa t})/(2\kappa)$. A flat curve gives $\kappa \cdot 4\% + \sigma^2(1 - e^{-2\kappa t})/(2\kappa)$. Vasicek's constant cannot match a whole curve, since it has too few parameters. In a simulation, the model gives back the ten-year zero of $e^{-0.4}$.

The drift is what a curve requires in order to be fitted exactly. The forward in the model is $f(t, T) = f(0, T) + \tfrac12\sigma^2\bigl(B(0, T)^2 - B(t, T)^2\bigr) + e^{-\kappa(T - t)}x_t$, with $B(t, T) = (1 - e^{-\kappa(T - t)})/\kappa$, $\mathrm{d}x = -\kappa x\,\mathrm{d}t + \sigma\,\mathrm{d}W$ and $x_0 = 0$. So it moves with a shock that falls off with maturity at the speed $\kappa$.

Simulate $r$ on small time steps with this drift, and sum it along each path. With the term above, the average discount matches the market to within noise. Without it, the model misses.

\iqwrong{Fitting the mean reversion from history and saying that fits the curve. In Hull-White, the speed comes from options and the drift comes from the curve.}
\end{solution}

\begin{solution}{iq:iv2:rates-and-fx:37}
The cap is worth about 59.9 basis points. The swaption is worth about 60, 52 and 43 at correlations of 1, 0.5 and 0. It is never worth more, because an option on an average is never worth more than the average of options. This is what links the cap and swaption markets: their gap reflects correlation.

The cap is $1.88 \times$ the caplet, $0.04(2\Phi(0.1) - 1)$. The swaption is $1.88E[(\tfrac12(F_1 + F_2) - K)^+]$. With $\rho = 1$, they are equal.

By convexity, the maximum of a sum is no more than the sum of the maxima. So a correlation below 1 lowers the average's variance, and that lowers the swaption.

\iqwrong{Calibrating a one-factor model to caps and to swaptions and expecting a fit. One factor means that forward rates are nearly perfectly correlated, so it overprices swaptions relative to caps.}
\end{solution}

\begin{solution}{iq:iv2:rates-and-fx:38}
The bullets are 104.45, 100.00 and 95.79. The calls are 3.86, 1.39 and 0.29. The callables are 100.60, 98.61 and 95.50. Their durations are 1.44, 2.60 and 3.76, against about 4.3 for the bullets. The callable's price barely rises as yields fall, which is negative convexity\index{negative convexity}. Its holder is short it.

At a low yield, the issuer calls, so the investor gets back par, and the bond acts like a one-year bond. At a high yield, the call is worthless, and the bond acts like the bullet. So duration shortens as yields fall, which is the opposite of what a holder wants.

A mortgage bond has the same shape, as borrowers refinance. A holder who hedges by selling more bonds as yields rise adds to the sell-off.

\iqwrong{Comparing yields without the option's value. A callable's extra yield is the premium for being short that call.}
\end{solution}

\begin{solution}{iq:iv2:rates-and-fx:39}
The fair price is about 97.6, and that is about 2.5 less than the lower of the two means. If yields rise, the long bond is cheapest. If they fall, the short bond is.

The two bonds both have a price of 100 at 6\%. At delivery, the price of the cheaper one is $\min(P_5(y), P_{20}(y))$. The short has an option, and it costs the long.

The long bond has much more duration. So the future's delta changes as the cheapest one changes, and it acts as if the long has sold a straddle to the short. The usual delta, using one bond, understates this.

\iqwrong{Pricing the future from today's cheapest bond and its carry only. That ignores a real option with real value when the bonds are near one another.}
\end{solution}

\begin{solution}{iq:iv2:rates-and-fx:40}
The three components explain about 85.3\%, 10.7\% and 2.6\%. The first has all maturities moving the same way, the second has short and long moving in opposite directions, and the third has the middle moving against the wings. These are level, slope and curvature\index{principal components}.

Take the covariance matrix, and find its eigenvectors. The first has weights of about 0.50, 0.59, 0.50 and 0.40. The second has $-0.74$, $-0.09$, 0.40, 0.53. The third has $-0.42$, 0.62, 0.19, $-0.63$.

So you can hedge the first two by holding two instruments that cancel them, and be left with the third. That is a butterfly weighted by these loadings, in place of a DV01-neutral one. However, the components are estimated from a sample, and their weights shift. So re-estimate them and watch them.

\iqwrong{Claiming the first component is a pure parallel shift. Its weights are close but different, which is why DV01 hedges leak.}
\end{solution}

\begin{solution}{iq:iv2:rates-and-fx:41}
The exact adjustment is about 3.41 basis points, and the first-order rule gives about 3.08. It arises because the payment date is out of line with the natural payment date.

Under the measure for the payment at $T + \tau$, $F$ is a martingale. To pay at $T$, you change measure with $P(T, T + \tau)^{-1} = 1 + \tau F_T$. So
\[
E^T[F_T] = \frac{E^{T+\tau}[F_T(1 + \tau F_T)]}{1 + \tau F_0} = F_0 + \frac{\tau F_0^2(e^{\sigma^2T} - 1)}{1 + \tau F_0}.
\]
With $\tau = 1$, $F_0 = 0.04$ and $\sigma^2T = 0.2$, this gives 3.41 basis points. The rule, $\tau F^2\sigma^2T/(1 + \tau F)$, gives 3.08.

Rate and discounting are linked. A high rate is paid early, at a time when it is worth more, and that is good for the holder.

\iqwrong{Using the forward with no adjustment. It is small here, but for a long-dated rate paid early, it is larger than the bid-ask.}
\end{solution}

\begin{solution}{iq:iv2:rates-and-fx:42}
The first-order rule gives about 17.6 basis points, and the simulation about 18.2. The adjustment raises the rate. It depends on the whole smile, as a simulation of the payoff is a function of the rate across all strikes.

Under the annuity measure, the swap rate $S$ is a martingale. To pay at the fixing date, you change measure: $E^T[S] = E^A[S/A(S)]\,/\,E^A[1/A(S)]$, where the denominator makes $E^T[1] = 1$. With a flat-yield form, $A(y) = (1 - (1 + y)^{-10})/y$. At 4\%, $A = 8.11$ and $A' = -40.4$. The rule is $F^2(e^{\sigma^2T} - 1)(-A'/A) = 0.0016 \times 0.2214 \times 4.98$.

Since $A$ falls as $S$ rises, $1/A$ is largest when the swap rate is high. So the payoff has a convex tilt. In the market, the integral goes across swaptions at all strikes, and so a smile that lifts the wings lifts this adjustment.

\iqwrong{Using the forward swap rate for the expected fixing. On long-dated CMS, the error can be many basis points.}
\end{solution}

\begin{solution}{iq:iv2:rates-and-fx:43}
The quanto\index{quanto} forward is about 100.60 at one year, above spot, and a simulation agrees. With a positive correlation, it would be below.

The index in dollars, $SX$, is a traded thing, where $X$ is the dollar value of a yen. So with zero rates, $SX$ has no drift under the dollar measure, and neither does $X$. By Ito, $d(SX)/(SX) = (\mu + \rho\sigma_S\sigma_X)\,dt + \ldots$, so the index's drift must be $\mu = -\rho\sigma_S\sigma_X$. That gives $100e^{0.3 \times 0.2 \times 0.1}$.

For the check, go through the yen measure. There, $S$ and $Y = 1/X$ are martingales, with a correlation of $+0.3$. The change to the dollar measure weights by $Y$, so $E^{\text{usd}}[S] = E^{\text{yen}}[SY]/E^{\text{yen}}[Y]$. The simulation gives 100.60.

The sign makes sense. A dollar investor who owns the real index gets hurt when the yen falls with the index up, as the correlation here says. The quanto is free of that, and it is worth more.

\iqwrong{Assuming the sign without saying what $X$ measures. The sign of the correlation flips when you quote the rate the other way, so define it first.}
\end{solution}

\begin{solution}{iq:iv2:rates-and-fx:44}
With no premium in the delta, the strike is $Fe^{\sigma^2T/2} \approx 1.1278$. When it includes the premium, it is $Fe^{-\sigma^2T/2} \approx 1.1166$. Pairs where the premium is paid in the foreign currency tend to use the second, as the premium itself is a position in that currency.

For a call's delta to be $-1$ times the put's, you set $\Phi(d_1) = \tfrac12$, so $d_1 = 0$. With the premium-adjusted form, the delta is $(K/F)\Phi(d_2)$, and the two legs sum to zero when $d_2 = 0$.

If you confuse them, then your 25-delta quote is at the wrong strike. A desk's smile fits badly, and its hedge is wrong.

\iqwrong{Treating the term at the money as a single point. There are several, and each market has its convention.}
\end{solution}

\begin{solution}{iq:iv2:rates-and-fx:45}
The forward is about 6.0\% just before two years, and about 3.6\% just after. It is 3.2\% just before three years, and 4.4\% just after. With log-linear discount factors, it is flat in each period at 5.0\%, 3.4\% and 4.8\%. You would hedge with the second, as the first's saw-teeth have no meaning.

For a linear zero rate, $f(t) = z(t) + tz'(t)$. Between one and two years, the slope is 1\%, so at two years, $4\% + 2 \times 1\% = 6\%$. Just after, the slope is $-0.2\%$, so $4 - 0.4 = 3.6\%$.

With piecewise-flat forwards, the rate is the log discount difference. Moving one input changes only a local segment. So its deltas stay in nearby buckets, while a smoother scheme, such as a cubic spline, can spread a bump to far places.

\iqwrong{Choosing a scheme by how a yield chart looks. The test is what happens to its forwards and to its risk.}
\end{solution}

\begin{solution}{iq:iv2:rates-and-fx:46}
The bond's price volatility is about 3.56\%. The call is worth about 1.42\%. The receiver swaption is about 1.44\%. They are nearly the same since each has the same exposure to one year of yield moves. The swaption has the larger exposure to its annuity.

$4.45 \times 0.8\% = 3.56\%$. The rule gives $0.3989 \times 3.56\% = 1.42\%$, and the swaption's is $4.5 \times 0.8\% \times 0.3989 = 1.44\%$. A call on a bond gains when yields fall, as does a receiver.

Their differences: a five-year bond rolls down to four years at expiry, so its forward duration at expiry is shorter. So use the forward duration. The forward bond's price is also funded through repo. So the two have different carry and specials risk.

\iqwrong{Using today's duration for a bond option, where the bond has aged by expiry. The one-year aging takes off about one point of duration.}
\end{solution}

\begin{solution}{iq:iv2:rates-and-fx:47}
The implied repo is 3\%. So you buy the bond, financing at 2.5\%, and sell the future, to lock in about 0.25\% in half a year.

The rate is $(99.5 + 2 - 100)/100 \times (1/0.5) = 3\%$. It is what you earn by holding the bond to deliver it. Financing it costs less. That is a cash-and-carry trade\index{basis trade}.

What can spoil it: a switch of the cheapest bond (you then hold the wrong bond); the short's delivery options, which you own and may not get value from; a change of the repo rate before delivery; and in a stress, the fall in the supply of repo, which makes leverage dangerous. These trades tend to be small per unit, and so they use high leverage.

\iqwrong{Calling this riskless. The basis has blown out before, when financing tightened, and the sizes grew fast.}
\end{solution}

\begin{solution}{iq:iv2:rates-and-fx:48}
The strike is about 4.08\%, on an annual-compounding basis. The swaption is worth about 1.338\% of notional. A simulation gives 1.335, with an error of 0.002.

A payer swaption is a put on a coupon bond with strike 1. In one factor, every bond price at the expiry falls in $r$. So there is a single $r^*$ at which the coupon bond equals 1. Each zero then has its own strike $X_i = P(1, T_i, r^*)$, and the option on the sum equals the sum of options on the zeros. This is Jamshidian's\index{Jamshidian decomposition} trick. Each zero's put has a closed form.

In two factors, the bond's value being below par is not one level. The set is a curve in the state space, and the order of the bond prices does not stay fixed, so it does not split up.

For the simulation, use the one-year-forward measure, in which the short rate's mean at expiry is the forward. Then discount with $P(0, 1)$.

\iqwrong{Using a flat-volatility Black number on a model price. A model's implied swaption volatility is an output, and it varies with the speed.}
\end{solution}

\begin{solution}{iq:iv2:rates-and-fx:49}
It replicates the vega, vanna and volga of the new option with three liquid ones. The new option's price is its flat price plus the weights times the smile cost of each pillar. It gives 10.41\% at 1.10 and 9.57\% at 1.25. The pillars come back exactly because the weights there are 1 and 0.

The weights solve a three-by-three system. That has a closed form:
\[
x_1(K) = \frac{\text{vega}(K)}{\text{vega}(K_1)}\frac{\ln(K_2/K)\ln(K_3/K)}{\ln(K_2/K_1)\ln(K_3/K_1)},
\]
and the same with the indices rotated. At $K = 1.25$, the weights are 0.33, $-0.96$ and 1.56.

The first-order formula gives good volatilities near the pillars. Outside the 25-delta range, the weights grow, and the smile can behave badly. At 1.30, it gives 9.97\%, rising fast from 9.57\%. So check for arbitrage before you use it in the wings, and use a model or a floor in the far tail.

\iqwrong{Calling vanna-volga a model. It is an interpolation rule with a hedging story. It has no dynamics, and it can be wrong on exotics.}
\end{solution}

\begin{solution}{iq:iv2:rates-and-fx:50}
The Europeans are worth about 1.14, 1.18, 0.95 and 0.54 as a percentage of notional, and the Bermudan about 1.585. The switch value over the best European is therefore about 0.40, or 34\%. With a speed of 0.3 and the same volatility, all the values fall: the best European is 0.76, the Bermudan 1.105 and the switch value 45\% of the best European. A faster speed raises the share, since the swap rates across dates become less correlated.

Set $x$ to be the driftless Gaussian state. On a grid in $x$, go back from year 4 to year 1. At each date, take the larger of the value of holding and the swap's exercise value. To step from year $t + 1$ to $t$, integrate under the $(t + 1)$-forward measure, whose transition is Gaussian with a known shift, and multiply by the discount factor $P(t, t + 1, x)$.

For the first check, the grid reprices a zero: it gives back $P(0, 3)$ to within $10^{-7}$. For the second, the first European agrees with Jamshidian's closed form.

A Bermudan is at least worth its best European. With a higher speed, shocks die off faster, so the later swap rates depend less on the earlier ones. Then the holder has more real choices, and the extra right is worth more compared with any one European. That is why Bermudan prices are read as a quote for the speed.

\iqwrong{Calibrating the model only to one set of Europeans and quoting a Bermudan. The switch value depends on the speed, and the Europeans do not pin that down.}
\end{solution}
